% ══════════════════════════════════════════════════════════════════
%  FICHE D'EXERCICES — Chapitre 5 : Les alcools et les phénols
% ══════════════════════════════════════════════════════════════════
\section{Exercices type Baccalauréat}
\textit{Données : $M(\ce{H})=1$, $M(\ce{C})=12$, $M(\ce{O})=16$ \si{\gram\per\mol}.}

\rubrique{Classes, nomenclature}

\exo{}\diff{1} Pour chacun des alcools suivants, donner le nom et la classe
(primaire, secondaire ou tertiaire) : a)~$\ce{CH3-CH2-CH2-OH}$ ;
b)~$\ce{CH3-CH(OH)-CH3}$ ; c)~$\ce{(CH3)3C-OH}$ ; d)~$\ce{CH3-CH(OH)-CH2-CH3}$.

\exo{}\diff{1} Écrire les formules semi-développées de tous les alcools de formule
brute $\ce{C4H10O}$, les nommer et préciser leur classe.

\rubrique{Oxydation ménagée}

\exo{}\diff{2} On oxyde de façon ménagée chacun des alcools suivants par une solution
acidifiée de dichromate de potassium en défaut puis en excès. Indiquer la nature du
(des) produit(s) et écrire les formules :
a)~propan-1-ol ; b)~propan-2-ol ; c)~2-méthylpropan-2-ol.

\exo{}\diff{2} Un alcool $A$ de formule $\ce{C3H8O}$ subit une oxydation ménagée
ménageant un produit $B$ qui rosit le réactif de Schiff et réduit la liqueur de Fehling.
\begin{enumerate}[label=\arabic*)]
\item Quelle est la classe de $A$ ? Justifier.
\item Identifier $A$ et $B$, les nommer, et écrire l'équation de l'oxydation
(le couple oxydant est $\ce{Cr2O7^2-}/\ce{Cr^3+}$).
\end{enumerate}

\exo{}\diff{2} Un alcool $C$ donne, par oxydation ménagée, un composé $D$ qui ne
réagit ni avec la liqueur de Fehling ni avec le réactif de Tollens, mais donne un
précipité jaune avec la 2,4-DNPH. Déterminer la classe de $C$ et la nature de $D$.

\rubrique{Déshydratation, estérification}

\exo{}\diff{1} Écrire l'équation de la déshydratation intramoléculaire du butan-2-ol
(en présence d'acide sulfurique à chaud). Donner les alcènes possibles et nommer le
produit majoritaire (règle de Zaïtsev).

\exo{}\diff{2} On fait réagir l'éthanol avec l'acide éthanoïque.
\begin{enumerate}[label=\arabic*)]
\item Écrire l'équation de la réaction d'estérification, nommer l'ester formé.
\item Quelles sont les caractéristiques de cette réaction (vitesse, limitation,
rôle du catalyseur) ?
\end{enumerate}

\rubrique{Détermination de formule}

\exo{}\diff{2} Un alcool saturé $A$ a une masse molaire $M = \SI{74}{\gram\per\mol}$.
Déterminer sa formule brute, écrire ses isomères alcools et préciser leur classe.

\exo{}\diff{3} La combustion complète de \SI{1.48}{\gram} d'un alcool saturé $A$
fournit \SI{3.52}{\gram} de dioxyde de carbone et \SI{1.80}{\gram} d'eau.
\begin{enumerate}[label=\arabic*)]
\item Déterminer la formule brute de $A$ (on montrera qu'il contient un atome d'oxygène).
\item Donner les isomères de $A$ et leur classe.
\item $A$ s'oxyde en un composé qui donne le test de Tollens positif. Identifier $A$
et écrire l'équation de son oxydation ménagée en excès d'oxydant.
\end{enumerate}

\rubrique{Problèmes de synthèse — type Baccalauréat}

\sujetbac[d'après Baccalauréat D, Cameroun]
Un composé organique $A$, de formule brute $\ce{C4H10O}$, présente les propriétés
suivantes : il réagit avec le sodium en dégageant du dihydrogène, et son oxydation
ménagée par le dichromate de potassium en milieu acide donne un composé $B$ qui
réduit la liqueur de Fehling.
\begin{enumerate}[label=\arabic*)]
\item À quelle famille appartient $A$ ? Quelle est sa classe ? Justifier.
\item Écrire les formules semi-développées possibles de $A$ et les nommer.
\item Écrire l'équation de la réaction de $A$ avec le sodium.
\item Écrire l'équation de l'oxydation ménagée de $A$ en $B$ (en excès d'oxydant,
le produit final est un acide carboxylique). Nommer le produit final.
\end{enumerate}

\sujetbac[Détermination d'un alcool par estérification]
On fait réagir \SI{6.0}{\gram} d'un alcool primaire saturé $A$ avec un excès d'acide
éthanoïque. On obtient \SI{8.8}{\gram} d'un ester, avec un rendement de \SI{100}{\percent}.
\begin{enumerate}[label=\arabic*)]
\item Établir la relation entre la masse molaire de $A$ et celle de l'ester.
\item Déterminer la masse molaire de $A$, sa formule brute et son nom.
\item Écrire l'équation de l'estérification et nommer l'ester.
\end{enumerate}
